% Generated draft: requires mathematical and source review.
\documentclass[11pt]{article}
\usepackage[margin=25mm]{geometry}
\usepackage{fontspec}
\setmainfont{Latin Modern Roman}
\usepackage{amsmath,amssymb,mathtools,unicode-math}
\setmathfont{Latin Modern Math}
\usepackage{enumitem,needspace,xurl,hyperref}
\setlength{\parindent}{0pt}
\setlength{\parskip}{5pt}
\setlength{\emergencystretch}{3em}
\newcommand{\N}{\mathbb{N}}
\newcommand{\Z}{\mathbb{Z}}
\newcommand{\Q}{\mathbb{Q}}
\newcommand{\R}{\mathbb{R}}
\newcommand{\C}{\mathbb{C}}
\newcommand{\F}{\mathbb{F}}
\newcommand{\A}{\mathbb{A}}
\newcommand{\PP}{\mathbb{P}}
\newcommand{\E}{\mathbb{E}}
\newcommand{\eps}{\varepsilon}
\newcommand{\dd}{\,\mathrm{d}}
\begin{document}
\section*{UnsolvedMath ID: 2302003}
\textbf{Source problem code:} catalogue entry
\subsection{Definitions and mathematical statement}
For a nonconstant entire function $h$, let
\[
M(r,h)=\max_{|z|=r}|h(z)|,\qquad
\operatorname{ord}(h)=\limsup_{r\to\infty}\frac{\log\log M(r,h)}{\log r}.
\]
Constants and polynomials have order zero in the convention used here. An entire function $\phi$ is an asymptotic function of an entire function $f$ along a path $\gamma:[0,\infty)\to\mathbb C$ if $\gamma$ is continuous, $|\gamma(t)|\to\infty$, and
\[
f(\gamma(t))-\phi(\gamma(t))\longrightarrow0\qquad(t\to\infty).
\]
Different asymptotic functions may use different paths, and distinct means not identical as entire functions.

\textbf{Source question.} Is it true that an entire function $f$ of order $\rho$ can have \emph{almost $2\rho$} distinct asymptotic functions $\phi$ of order less than $1/2$? The extracted word is almost, not at most. As written, it suggests a near-attainment question but gives no tolerance or quantifier. If the intended wording was at most, the question instead asks for an upper bound on the cardinality of the family. These are different targets, so neither is silently selected. The opening phrase growing slowly compared with $f$ also lacks a quantitative definition beyond the later order restriction.

The source reports a positive result due to Denjoy when the asymptotic paths are straight lines, and a polynomial-asymptotic-function case as a consequence of Ahlfors' theorem for asymptotic values. It also states that if $\psi=\phi_1-\phi_2$ has minimum modulus tending to zero, then its lower order is at least $1/2$ with a mean-type qualification. Here minimum modulus is $\min_{|z|=r}|\psi(z)|$. The exact lower-order/type assertion and its reference are incomplete in the extraction. These are retained as background reports rather than used to decide which version of the main question was intended.
\subsection{Short English statement}
Relate the number of small-order entire asymptotic functions of $f$ to the order of $f$. The source’s almost wording leaves an essential ambiguity between a bound and a near-attainment problem.
\subsection{Sources}
\begin{itemize}
\item Hayman, \emph{Research Problems in Function Theory} (2018), Problem 2.3, \url{https://arxiv.org/abs/1809.07200}; source-TeX extraction, accessed 2026-07-29. The 2018 update was retained for status triage.
\item Denjoy, straight-line-path result; Ahlfors, theorem on asymptotic values; an incomplete see-Hayman reference for the minimum-modulus/lower-order observation. The supplied text omits the corresponding bibliographic entries and theorem locations.
\item Local registry \texttt{data/unsolvedmath/problems.json}, permanent ID $2302003$, source code catalogue entry, \emph{Research Problems in Function Theory --- Problem 2.3}. The record labels this open and OPEN-TRIAGE, reporting no recent resolution located, with NEEDS\_REVIEW source-status and rights qualifications.
\end{itemize}
\end{document}
